\documentclass{syllabus}

\begin{document}

\thispagestyle{fancy}
\coursecode{PHY 3970}
\date{Spring 2026}
\author{}
\title{\coursecode: {Portfolio Seminar C}}
\maketitle

\meeting{ Tuesday: & 8:00 AM - 9:40 AM, DSC 054 }
\officehours{TR: & 11:30 AM - 12:30 PM}
\prerequisites{PHY 3960: & (Portfolio B)}


\section*{General Information}

\infotableport





\section*{Course overview}

\textbf{Catalog description:} The third course in the Physics Portfolio Seminar presents students with more opportunities to approach scenario-style problems using physics principles. Emphasis is on group work and the informal and formal presentations of solutions to increasingly complex open-ended questions. Students will also continue to build and reflect upon their physics portfolios.


The purpose of this course is to get hands-on experience designing, calibrating, and deploying a laboratory apparatus, to include
\begin{enumerate}
    \item Identifying basic principles of measurement,
    \item Designing and building supporting laboratory apparatus to make use of available sensors,
    \item Modeling and calibrating the measurement,
    \item Performing measurements and analysing the resulting data.
\end{enumerate}




\subsection*{Student learning outcomes} 
\noindent

Upon successful completion of this course, students will be able to:
\begin{enumerate}[label=\textcolor{CarthageDark}{\bfseries\arabic*.}]
    \item Make design choices in light of physical principles and measurement goals,
    \item Use basic fabrication skills to create an apparatus,
    \item Make quantitative models of the expected behavior of an instrument,
    \item Adjust the interpretation of results based on calibration measurements,
    \item Interpret sensor data to draw physical conclusions.
\end{enumerate}



\section*{Course components}

\paragraph{Design:} Students will determine a method of measurement and simulate the expected behavior of the system, making justified choices for design parameters based on a quantitative model.

\paragraph{Build:} Students will use basic fabrication tools and/or acquire equipment and assemble their apparatus. A plan for and report from the build phase will be required.

\paragraph{Calibrate and measure:} At least two experiments will be performed using the designed apparatus. A combined report on these experiments will be written and presented orally at the end of the course.



\section*{Grading Scheme}
\subsection*{Component Weights}

\noindent\begin{tabular}{lr}
\textbf{Component} & \textbf{Weight}       \\ \hline
\textbf{Design}      & 35\%  \\
\textbf{Build}           & 30\%       \\
\textbf{Measure}            & 30\%     
\end{tabular}%



\lettergrades{a) demonstrated effort in class, b) willingness to seek help outside of class.}


\tipsforsuccess{problem}

\section*{Policies}
\subsection{General course policies}
Expect a response within one business day for emails sent to the instructor. To ensure a prompt response, format the subject line as: [\coursecode] *insert subject here*

All required material will be covered in lecture, but can also be found in the parallel readings, so in the case of absence students should consult the scheduled readings. Students found sleeping will be brought to the attention of the class for public ridicule and no effort made to wake them up.

Students are expected to arrive and leave on time. I will attempt to arrive five minutes early and be available for at least five minutes after class.

\latework{None}

\collegepolicies[Other college policies]

\end{document}
